\documentclass[11pt]{article}
\usepackage[margin=1in]{geometry}
\usepackage[T1]{fontenc}
\usepackage{lmodern,amsmath,amssymb,amsthm,mathtools,bm,booktabs,enumitem}
\usepackage[round,authoryear]{natbib}
\usepackage[colorlinks=true,linkcolor=blue,citecolor=blue,urlcolor=blue]{hyperref}
\setlength{\parindent}{0pt}
\setlength{\parskip}{5pt}
\setlength{\emergencystretch}{2em}
\setlist{itemsep=2pt,topsep=3pt,leftmargin=*}
\newcommand{\E}{\mathbb E}
\newcommand{\Prb}{\mathbb P}
\newcommand{\Var}{\operatorname{Var}}
\newcommand{\Cov}{\operatorname{Cov}}
\newcommand{\tr}{\operatorname{tr}}
\newcommand{\diag}{\operatorname{diag}}
\newcommand{\op}{\mathrm{op}}
\newcommand{\ind}{\mathbf 1}
\newcommand{\calA}{\mathcal A}
\newcommand{\calB}{\mathcal B}
\newcommand{\calT}{\mathcal T}
\newcommand{\calP}{\mathcal P}
\newcommand{\calE}{\mathcal E}
\newtheorem{proposition}{Proposition}
\newtheorem{lemma}[proposition]{Lemma}
\newtheorem{assumption}{Assumption}
\title{RoCE-K\\[3pt]\large Method design for weak source bias and increasing site counts}
\author{Sinian Zhang}
\date{October 2026\\\small Design version 1, 1 October 2026}
\begin{document}
\maketitle
\begin{center}
\fbox{\parbox{0.93\linewidth}{\textbf{Research design for discussion.}
This document specifies a candidate method and conditional mathematical reductions.
It does not claim a completed high-dimensional RoCE-K theorem, a validated final
confidence procedure, or a new optimality result. The current RoCE manuscript
and its production defaults are separate.}}
\end{center}
\section{Design decisions before further simulation}

The objective is inference for a target average treatment effect when some
source arms have small, hard-to-detect departures and the source count
$K=K_n$ can increase. The primary requirements are coverage over local-bias
sequences, useful interval length when borrowing is informative, and a
computable procedure with explicit assumptions. Exact classification of every
source is not a requirement. Point-estimation risk is a separate objective;
coverage alone will not be used to assert an RMSE guarantee.

\begin{center}
\begin{tabular}{@{}p{.24\linewidth}p{.69\linewidth}@{}}
\toprule
Component & Proposed first design\\
\midrule
Candidate scores & One common target outcome prediction per arm, with
source-specific calibrated merged weights.\\
Aggregation & A joint two-arm working-covariance GLS contrast; arm weights
need not be equal.\\
Bias protection & Bound the aggregate bias through a noise-corrected quadratic
statistic, valid-count assumptions and explicit valid-score bias budgets.\\
Covariance & An independent pilot supplies a covariance upper confidence bound,
not an empirical covariance silently treated as known.\\
Safeguard & Combine with a separately valid target reference interval, which
may use the full target sample.\\
First proof design & Independent training, covariance-pilot and evaluation
samples. Efficient full cross-fitting is a later extension.\\
\bottomrule
\end{tabular}
\end{center}

The proposed change relative to the previous Gaussian reference is to estimate
bias dispersion with an unbiased quadratic statistic. Its conditional
variance can be bounded using second moments of independent patient blocks.
This avoids requiring a high-dimensional Gaussian approximation for that
statistic. The first calibration below is deliberately conservative; making
its constants useful at $n=1000$ is an explicit design gate.

\paragraph{Relation to existing work.}
The voting construction of \citet{guo2018} and the selection-robust searching
and sampling approach of \citet{guo2023} motivate inference that tolerates
uncertain validity. RIFL already studies robust federated inference
\citep{rifl2025}. Our source-assisted scores share target observations and have
two potentially different arm-validity sets; their covariance cannot be
replaced by independent site-estimate variances. Adaptive Gaussian interval
lengths, including a fourth-root term, are already studied by
\citet{wang2026}. Neither that rate nor the use of a quadratic statistic is
claimed as an innovation here. The proposed combination of calibrated causal
scores, shared-target dependence, estimated covariance, arm-specific validity
and increasing $K$ needs its own proof and comparison.

\section{Statistical contract}

There is a target site $t$ and sources $s_1,\ldots,s_K$. The primary outcome
is binary, or more generally bounded in $[0,1]$. Write
\[
 \mu_t^a=\E_t Y(a),\qquad \tau_t=\mu_t^1-\mu_t^0,
 \qquad a\in\{1,0\}.
\]
Identification at the target requires consistency, overlap and adequate
confounding adjustment. A source arm is structurally valid under the declared
transport and working-model conditions. Let $\mathcal V_a^\star$ denote this
unknown set. Specify lower bounds
\[
 |\mathcal V_a^\star|\ge g_a,\qquad 0\le g_a\le K,
\]
before examining agreement among estimates. The two sets may differ. Other
source-arm departures may be local, fixed, of either sign, or cancel in the
treatment contrast. There is no minimum separation assumption.

The target identifies $\tau_t$ even if no source is valid. The bounds $g_a$
are additional borrowing assumptions. They cannot be obtained by counting
apparently compatible sources and then treated as certain. If neither arm has
a guaranteed valid source, the prespecified default is the target reference
procedure. A sensitivity family over $(g_1,g_0)$ may be reported, with each
interval labeled by its assumption; the narrowest member must not be selected
as an assumption-free result.

\paragraph{Source counts and sample sizes.}
State $n_t$, each $n_{s_j}$, the nuisance dimension $p$, and sparsity parameters
separately. Positive site fractions of a fixed total sample cannot be assumed
uniformly when $K\to\infty$. An experiment that varies finite $K$ while holding
$n_j=1000$ fixed is a finite-sample study, not a growing-$K$ theorem.

\section{Sample roles and the candidate program}

\subsection{Honest first construction}
At each site, use an index-only random split into training $\calT$, pilot
$\calP$ and final evaluation $\calE$. All learned preprocessing, nuisance
fitting and nuisance tuning for these scores use $\calT$ only. Feature columns
and coding must agree across sites and between outcome and weighting models.
The fitted transformation is applied to both held-out samples. At
$1000$/site, $500/250/250$ is a diagnostic starting allocation, not an optimal
allocation claim. Treatment- or outcome-stratified splitting does not supply
the simple conditional-iid argument used below without an additional proof.

Conditional on training and its fitting randomness, patient observations in
the pilot and final samples are independent across roles and sites. Pilot
quantities may determine the working covariance and optional screening, but
the evaluation observations may not determine those choices. These are three
\emph{sample roles}, not the existing three-layer nested cross-fitting algorithm.

\subsection{Prefer a common outcome score to a translated source score}
Fit a common outcome regression $\widehat m_c^a$ at the target by an
inverse-propensity weighted outcome loss. For a logistic working model, the
population loss is
\[
 \E_t\!\left[\frac{\ind(A=a)}{\pi_t^a(X)}
 \{\log(1+e^{\beta^T\phi(X)})-Y\beta^T\phi(X)\}\right].
\]
Its empirical version uses training-only propensity fits, a declared penalty
policy, and fold-summed loss blocks where needed. The propensity weight here
is $1/\pi_t^a$, not the odds weighting used by the target-anchor OR. Retain
the separately calibrated target AIPW anchor $\widehat\psi_t^a$.

Each source retains its own calibrated merged weight $\widehat q_j^a$.
For evaluation patient $i$, define
\begin{align}
 R_{ji}^a&=\ind(A_{ji}=a)\widehat q_j^a(X_{ji})
                \{Y_{ji}-\widehat m_c^a(X_{ji})\},\label{eq:residual}\\
 Z_j^a&=\mathbb P_{t,\calE}\widehat m_c^a
                 +\mathbb P_{s_j,\calE}R_j^a,\qquad
 Z_t^a=\mathbb P_{t,\calE}\widehat\psi_t^a.\label{eq:candidate}
\end{align}
The same fitted $\widehat m_c^a$ appears in the target prediction and every
source residual. Source averages divide by the entire site sample size,
including the treatment indicator in~\eqref{eq:residual}.

This choice removes the additional term
$(\mathbb P_{t,\calT}-\mathbb P_{t,\calE})(\widehat m_j^a-\widehat m_c^a)$
introduced by translating source-specific predictions. It is selected for
its structure, not because the previous fitted diagnostics showed uniformly
smaller bias; they did not. The translated-source and original source-specific
scores remain diagnostic comparators. A dedicated implementation may omit
final source OR fits, but must first reproduce the retained weight fits.

\subsection{The robustness branches must be explicit}
For a valid source, with $m_\beta=\operatorname{expit}(\beta^T\phi)$ and
$q_\gamma=\exp(-\gamma^T\phi)$ away from truncation boundaries, the population
score $S_j=\E_t m_\beta+\E_j\{\ind(A=a)q_\gamma(Y-m_\beta)\}$ has gradients
\begin{align*}
 \partial_\beta S_j&=\E_t(\dot m_\beta\phi)
              -\E_j\{\ind(A=a)q_\gamma\dot m_\beta\phi\},\\
 \partial_\gamma S_j&=-\E_j\{\ind(A=a)q_\gamma(Y-m_\beta)\phi\}.
\end{align*}
Both vanish under either of these population branches:
\begin{enumerate}
\item The source joint weight and target propensity are correct. The common
IPW outcome limit is the full-target projection, giving the required residual
moment even when the outcome working model is misspecified.
\item The common OR is correct, and source weight calibration balances its
limiting derivative-weighted features. The initial OR used by calibration
must have the correct limit in this branch.
\end{enumerate}
Correct source weights alone do not justify the first branch if both the
target propensity and common OR are misspecified. Active population clipping,
estimated IPW weights and CV-selected penalties require their own treatment.
We do not assume that this candidate inherits every robustness configuration
of the source-specific-OR procedure.

\section{Joint score representation and two uncertainty inputs}

Stack candidates in arm-major order,
\[
 Z=(Z_t^1,Z_1^1,\ldots,Z_K^1,Z_t^0,Z_1^0,\ldots,Z_K^0)^T,
 \qquad H=\begin{pmatrix}\mathbf1_{K+1}&0\\0&\mathbf1_{K+1}\end{pmatrix}.
\]
Let $\mu=(\mu_t^1,\mu_t^0)^T$ and $\ell=(1,-1)^T$. Conditional on training,
\begin{equation}
 \E Z=H\mu+r,\qquad \Cov(Z)=\Sigma_E.\label{eq:score-model}
\end{equation}
The vector $r$ includes nuisance drift at valid coordinates as well as
transport departures. Structural validity does not make a fitted candidate
conditionally unbiased.

The target supplies the four-dimensional patient score
$G_t=(\widehat\psi_t^1,\widehat\psi_t^0,
\widehat m_c^1,\widehat m_c^0)^T$; source $j$ supplies
$G_j=(R_j^1,R_j^0)^T$. For a known linear assembly map $L$,
\begin{equation}
 Z=L(\overline G_t^T,\overline G_1^T,\ldots,\overline G_K^T)^T,
 \quad \Sigma_E=L\diag\{\Omega_b/n_b^E\}_{b=t,1,\ldots,K}L^T,
 \quad \Omega_b=\Cov(G_{bi}\mid\calT).\label{eq:blocks}
\end{equation}
Within-person cross-arm covariance is retained. Independence of patients
does not imply independence of two arm estimates. Every matrix and covariance
normalization below refers to final mean sample sizes $n_b^E$.

\subsection{Input A: valid-score bias envelopes}
Require an event $\mathcal E_B$, with failure probability at most $\delta_B$,
on which
\begin{equation}
 |r_i|\le b_i\quad\text{for all truly valid coordinates, including the anchors}.
 \label{eq:bias-event}
\end{equation}
The supplied $b_i$ is a bound that applies \emph{if} coordinate $i$ is valid;
it does not restrict an invalid coordinate. These bounds must be justified by
the nuisance program, including truncation bias. An oracle population
integration error measured in a simulation is not an implementable budget.

A prospective sparse-model route is a simultaneous bound of order
$s_{\rm eff}\log(pK/\delta_B)/n^T_{\min}$ under proved calibration, curvature
and penalty conditions. Its constants and event probability are not supplied
by that rate notation. Obtaining a usable budget, or proving its omission is
uniformly negligible on the intended interval scale, is the principal open
high-dimensional step. The zero-budget option is reserved for justified
unbiased-score cases or explicitly labeled oracle diagnostics.

\subsection{Input B: an upper covariance certificate}
Use the independent pilot to construct a positive definite $W$ such that
\begin{equation}
 \Prb(\Sigma_E\preceq W\mid\calT)\ge1-\delta_\Sigma.
 \label{eq:cov-event}
\end{equation}
The proposal needs an upper bound, not a lower Loewner bound. One feasible
but potentially loose construction is given in Appendix~\ref{app:covariance}.
It estimates at most $4\times4$ target and $2\times2$ source blocks and then
assembles~\eqref{eq:blocks}. A tighter pilot certificate can replace it only
after its simultaneous coverage is established.

The pilot working covariance $W$ and the evaluation covariance correction
$\widehat\Sigma_E$ have different roles. For the latter, use the unbiased
within-site covariance
\begin{equation}
 S_b^E=\frac{1}{n_b^E-1}\sum_{i\in\calE_b}
 (G_{bi}-\overline G_b)(G_{bi}-\overline G_b)^T,
 \quad\widehat\Sigma_E=L\diag\{S_b^E/n_b^E\}_bL^T.\label{eq:eval-cov}
\end{equation}
The divisor $n_b^E-1$ is essential for the exact identity below. Covariance
and means may use the same evaluation records in~\eqref{eq:eval-cov}; their
dependence is handled by the U-statistic calculation. They must not be used
to refit $W$ or choose the aggregation coefficients.

\section{Joint aggregation and a bias-energy upper bound}

\subsection{Work in a pilot-chosen covariance metric}
Let $\calA$ contain the two target coordinates and any source coordinates known
valid by assumption, such as all control sources if $g_0=K$. Do not enlarge it
because observed estimates happen to agree. Write
\begin{align}
 a_W&=W^{-1}H(H^TW^{-1}H)^{-1}\ell,& v_W&=a_W^TWa_W,\\
 P_W&=W^{-1}-W^{-1}H(H^TW^{-1}H)^{-1}H^TW^{-1}.\nonumber
\end{align}
Embed the reference residual precision in the full coordinate space:
\begin{align}
 R_{\calA}&=\operatorname{embed}_{\calA}\!\left[
 W_{\calA}^{-1}-W_{\calA}^{-1}H_{\calA}
 (H_{\calA}^TW_{\calA}^{-1}H_{\calA})^{-1}
 H_{\calA}^TW_{\calA}^{-1}\right],\\
 D_W&=P_W-R_{\calA},\qquad \nu=2(K+1)-|\calA|.
\end{align}
Then $D_WH=0$, $D_WWD_W=D_W$, and $D_W$ is positive semidefinite of rank
$\nu$. The center $T_W=a_W^TZ$ uses separate arm weights optimized jointly
for TATE. $v_W$ is a \emph{variance upper proxy}, not the assertion
$\Var(T_W)=v_W$.

Define the conditional bias energy
\begin{equation}
 \calB=r^TD_Wr,\qquad
 \widehat{\calB}=Z^TD_WZ-\tr(D_W\widehat\Sigma_E).\label{eq:debiased-energy}
\end{equation}
The correction subtracts estimated sampling dispersion. In particular,
negative $\widehat{\calB}$ values are legitimate; the statistic itself is not
clipped to zero before applying its calibrated upper-bound rule.

\begin{proposition}[Conditional unbiasedness and variance bound]\label{prop:ustat}
Condition on training and pilot, so $W$ and $D_W$ are fixed. Suppose the
evaluation patient blocks are independent, iid within each site, have finite
second moments, and $n_b^E\ge2$. Then
\[
 \E\widehat{\calB}=\calB.
\]
On $\Sigma_E\preceq W$, with $\kappa_E=\max_b n_b^E/(n_b^E-1)$,
\begin{equation}
 \Var(\widehat{\calB})\le4\calB+2\kappa_E\nu.\label{eq:energy-var}
\end{equation}
No Gaussian approximation around invalid-source means and no separation of
valid and invalid sources is used in this conditional statement.
\end{proposition}

For $c_D=(1-\alpha_D)/\alpha_D$, define
\begin{equation}
 U_{\calB}=\max\!\left\{0,\widehat{\calB}+2c_D+
 \sqrt{\max\{0,4c_D\widehat{\calB}+4c_D^2+2c_D\kappa_E\nu\}}\right\}.
 \label{eq:energy-upper}
\end{equation}
Cantelli's inequality and~\eqref{eq:energy-var} give
$\Prb(\calB\le U_{\calB}\mid\calT,\calP)\ge1-\alpha_D$ on the covariance
event. If $\nu=0$, set $U_{\calB}=0$ and omit the dispersion step; its error
allocation may be reassigned by a rule fixed before observing data.
Appendix~\ref{app:quadratic} gives the algebra and explains the zero
within-person diagonal. This is a second-moment reference calibration;
its constants are not claimed to be practically optimal.

\subsection{Convert energy to a TATE bias allowance}
Let $V$ contain $\calA$ and the prescribed valid-source counts in each arm.
Choose any comparison coefficients $d_V$ supported on $V$ with
$H^Td_V=\ell$, and set
\begin{equation}
 \widetilde d_V=(I-R_{\calA}W)d_V,\quad
 B_V=\sum_{i\in V}|\widetilde d_{V,i}|b_i,\quad
 C_V=\widetilde d_V^TW\widetilde d_V-v_W.\label{eq:comparison-pair}
\end{equation}
For the true valid subset, $C_V\ge0$ and
\begin{equation}
 |a_W^Tr|\le B_V+\sqrt{C_V\calB}.\label{eq:bias-geometry}
\end{equation}
Thus a safe count-based bound maximizes the \emph{paired} expression in
\eqref{eq:bias-geometry} over admissible $V$. A separate reference comparison
supported on $\calA$ gives $(B_{\calA},C_{\calA})$. The proposed implementation
uses
\begin{equation}
 R_{\rm bias}(u)=\min\{B_{\calA}+\sqrt{C_{\calA}u},\;
                         \overline B_{\rm count}+\sqrt{\overline C_{\rm count}u}\},
 \label{eq:bias-radius}
\end{equation}
where the count pair must upper-bound the same admissible comparison family.
Independent selection of unrelated small bias and variance constants is not
valid. Exact subset enumeration is a small-$K$ check, not the large-$K$ plan.
Appendix~\ref{app:count-bound} specifies a conservative structural alternative.

\section{Reported interval, point estimate and coverage target}

\paragraph{Operational sequence.}
\begin{enumerate}
\item Declare $(g_1,g_0)$, the error budget, score bounds and nuisance policy;
split observations by index into training, pilot and evaluation roles.
\item Fit the common outcome models, source merged weights and core anchors
on training data. Construct the separately justified full-target reference.
\item Form the pilot covariance upper certificate $W$ and fix the joint GLS
coefficients and the reduced precision $D_W$.
\item Obtain evaluation score sums and cross-moments. Compute $Z$,
$\widehat\Sigma_E$, $\widehat{\calB}$ and $U_{\calB}$.
\item Compute the paired bias/variance comparison bounds and form $I_D$.
\item Apply the declared intersection/fallback rule, project the reported
point into the final interval, and retain every diagnostic component.
\end{enumerate}

Chebyshev gives the noise radius $r_C=\sqrt{v_W/\alpha_N}$. If valid score
ranges from Appendix~\ref{app:covariance} are available, let $L_b$ denote
the assembly matrix for site $b$ and put
\begin{align*}
 M_W=\max_b\frac{2}{n_b^E}\sum_l |(L_b^Ta_W)_l|L_{b,l},\qquad
 t_N&=\log(2/\alpha_N),\\
 r_B&=\frac{M_Wt_N}{3}+\sqrt{2v_Wt_N+(M_Wt_N/3)^2}.
\end{align*}
Here $M_W$ bounds each centered patient contribution to $T_W$.
The two-sided Bernstein inequality justifies $r_B$. Use
$r_N=\min(r_C,r_B)$ when these ranges are certified, and $r_N=r_C$ otherwise.
This choice is fixed conditional on training and pilot, so it requires no
additional error allocation. Define
\begin{equation}
 I_D=\left[T_W-\{r_N+R_{\rm bias}(U_{\calB})\},
           T_W+\{r_N+R_{\rm bias}(U_{\calB})\}\right].
 \label{eq:core-interval}
\end{equation}
This is a conservative validity reference. A scalar normal critical value
requires a separate uniform scalar-score approximation and its error
allowance. A sharper one-sided calibration of the zero-diagonal quadratic
statistic is a priority refinement; favorable simulated coverage does not
justify replacing the critical value.

Let $I_R$ be a separately valid target reference at failure level $\alpha_R$.
It may use \emph{all} target observations with its own valid cross-fitted
inference. Independence from $I_D$ is unnecessary. This avoids making a
quarter-sample target interval the only safeguard. Distinguish this external
reference interval from the internal comparison coefficients used in
\eqref{eq:comparison-pair}.

Return $I_D\cap I_R$ when it is nonempty. Otherwise return the shorter of the
two intervals, with a fixed tie rule and a recorded fallback flag. The output
length is pointwise at most either component length. For binary outcomes,
project its endpoints onto $[-1,1]$; this cannot remove the true parameter on
a coverage event. Report the projection of $T_W$ onto the final interval as
the point estimate, while retaining $T_W$ as a diagnostic. The reported point
therefore lies inside its interval. Its RMSE still needs an independent
analysis; it does not inherit an oracle MSE claim from GLS.

\begin{proposition}[Conditional reduction to feasible uncertainty inputs]
Suppose the validity-count assumption holds, the bias event has failure
probability at most $\delta_B$, the covariance event has failure probability
at most $\delta_\Sigma$, and $I_R$ has coverage at least
$1-\alpha_R-\epsilon_{R,n}$. For the independent-sample construction above,
\[
 \Prb(\tau_t\in I)\ge
 1-\delta_B-\delta_\Sigma-\alpha_N-\alpha_D-\alpha_R-\epsilon_{R,n}.
\]
This proposition is conditional on justified inputs $b$ and $W$; it is not a
completed theorem for fitted high-dimensional calibrated nuisances.
\end{proposition}
\begin{proof}
On the bias and covariance events, apply Chebyshev or the certified Bernstein
bound to the fixed-pilot linear contrast, Cantelli to~\eqref{eq:debiased-energy},
and the deterministic bias
bound. A union bound gives coverage of $I_D$. If both component intervals
contain the truth, their intersection is nonempty and contains it. A second
union bound with the reference failure event proves the result. No independence
between the two final intervals is required.
\end{proof}

For a prespecified borrowing analysis, one possible allocation is
$\delta_B=\delta_\Sigma=\alpha/[10(K+1)]$,
$\alpha_R=\alpha_D=\alpha/3$, and
$\alpha_N=\alpha-\alpha_R-\alpha_D-\delta_B-\delta_\Sigma$.
These choices are a design starting point. The reference is then more
conservative than a nominal $1-\alpha$ target interval; both references must
be shown in comparisons. In the prespecified no-borrowing case, use the
original target procedure at level $1-\alpha$.

\section{What increasing \texorpdfstring{$K$}{K} should mean for this method}

The conditional moment result does not introduce a growing-dimensional
Gaussian-approximation error. Coverage nevertheless needs simultaneous bias
and covariance certificates for all participating sites. Increasing $K$
still enters these certificates, numerical stability and interval length.

At an ideal all-valid conditional law with $r=0$,~\eqref{eq:energy-var} gives
$\E(\widehat{\calB})_+=O(\sqrt\nu)$. The inversion obeys
\[
 U_{\calB}\le2(\widehat{\calB})_++8c_D+\sqrt{2c_D\kappa_E\nu},
 \qquad \E\sqrt{U_{\calB}}=O(\nu^{1/4}).
\]
If, additionally, comparable site sizes give $v_W=O(1/(nK))$,
$\overline C_{\rm count}=O(1/(nK))$, $b=0$, and $\nu=O(K)$, then
\begin{equation}
 \E|I_D|=O\{n^{-1/2}K^{-1/4}\}.\label{eq:ideal-length}
\end{equation}
This is a conditional, ideal-score rate calculation. It is not an established
rate for the proposed fitted estimator. A common target variance adds its
own floor; a covariance upper bound that introduces a substantial artificial
common variance can also prevent this gain. An unconditional expected-length
claim additionally needs control of certificate failures and integrability
of the fallback length. Fixed nonvanishing failure allocations cannot simply
be ignored in such an argument.

For fitted candidates, a desired budget of order
$s_{\rm eff}\log(pK)/n$ suggests the sufficient scale requirement
\[
 \frac{s_{\rm eff}\log(pK)K^{1/4}}{\sqrt n}\longrightarrow0
\]
to make that budget negligible relative to~\eqref{eq:ideal-length}.
This rate must be proved for the actual training program. It is not supplied
automatically by minimum-CV-loss tuning. Covariance inflation, coefficient
norms and the shared-target variance floor need separate bounds.

Existing lower-bound calculations in a randomized binary submodel show why
uniform weak-bias inference cannot generally have the oracle-valid-set
$1/\sqrt{nK}$ width. The analogous Gaussian adaptation phenomenon is already
in the literature \citep{wang2026}. The objective is useful, honestly calibrated
borrowing within these limits, not perfect weak-source screening.

\section{Communication and computation}

Conditional on an agreed training-only feature map, the one-round version
can batch these messages:
\begin{enumerate}
\item Target sends the common OR coefficients, the block-specific initial
OR coefficients and derivative-weighted target moments needed for source
weight calibration, plus split and feature identifiers.
\item Each source fits its merged weights locally and returns, separately for
pilot and evaluation samples, its two residual sums, three distinct second
cross-moments, sample counts and justified score-range bounds.
\item Target assembles the common four-score block, $W$, the joint contrast,
the corrected quadratic statistic, and the final interval locally.
\end{enumerate}
Patient records remain local in this protocol. Learned pooled preprocessing
may require a separate agreement step and is not automatically included in
the estimation-round count. Source-local initial ORs retain the two-round
dependency for target derivative summaries. Changing sample roles does not
by itself change the communication depth.

The shared-target representation is a low-rank common block plus private
$2\times2$ blocks. Use block solves and Woodbury identities for GLS and the
quadratic trace calculation. Sorting private variance bounds removes subset
enumeration in the structural upper bound. A claim of linear or near-linear
total inference cost must include the bias-budget projection step; the
current dense reference is not claimed to have that complexity. Proximal
Newton and coordinate descent should remain explicit nuisance-solver options,
with matched tolerances and result-parity checks.

\section{Extensions after the core design is fixed}

\paragraph{Optional pilot-certified exclusion.}
Strongly discrepant source arms may be excluded on the pilot only when a
simultaneous confidence bound for their difference from the target exceeds
the combined valid-score budgets. On the certification event, no valid arm
is removed, so $g_a$ remains the original lower bound. An uncertain arm is
retained. A certification failure allowance must be added to the error budget.
If fewer than $g_a$ arms remain, return the prespecified reference and record
the contradiction; do not lower $g_a$ to rescue borrowing. Screening is off
in the first implementation.
If enabled later, restrict $Z,H,W$ to the retained arm coordinates and use
the actual retained dimension in $\nu$; retaining a source in only one arm
does not change the whole-site denominator of its residual mean.

\paragraph{Cross-fitting and reuse of data.}
The first proof concerns the independent three-role construction. A safe
intermediate extension is to construct intervals from prespecified rotations
with error budgets summing to $\alpha$, then intersect them with a declared
empty-set fallback. Dependence between rotations does not invalidate the
union bound, but its length cost must be assessed. Efficient averaging of
overlapping-fold scores requires an additional argument: conditioning on
all training sets does not make their evaluations jointly independent.
The current paper's two-layer estimator and its three-layer diagnostic remain
separate from this proposed extension.

\paragraph{Failures and reporting.}
Record fit failures, nonpositive working covariance, extreme budgets and
empty intersections. Use a justified target reference, or a parameter-space
interval when no usable reference is available. Never drop repetitions or
infer new validity counts from favorable results. Save the raw GLS point,
reported point, both interval components, noise and bias radii, corrected
dispersion, covariance certificate and all declared budgets.

\section{Development gates and experiment design}

\begin{enumerate}
\item \textbf{Freeze the statistical contract.} Review the common-outcome
score, $(g_1,g_0)$ assumptions, target reference, interval target and penalty
policy. The zero-budget and known-covariance versions are diagnostic modes.
\item \textbf{Check the conditional algebra.} Reconstruct the quadratic
statistic from moments and from patient pairs; enumerate small non-Gaussian
laws; check the variance bound, bias geometry, and arm covariance. These are
identity checks, not an MC coverage claim.
\item \textbf{Make the uncertainty inputs useful.} Derive uniform valid-score
budgets for the calibrated fits and a covariance certificate that is not
vacuous at practical sample sizes. Sharpen the elementary calibration before
claiming efficiency from a conservative reference construction.
\item \textbf{Validate a complete fitted procedure.} Regenerate and refit
every repetition. Start at small $p$, then $p=100,200$, using local departures
such as $h/\sqrt n$, both signs, different arm-validity sets and all-valid
cases. Vary $n$ jointly with $K$; keep the $1000$/site panel as a practical
slice rather than the only growth experiment.
Hold the invalid-source fraction fixed in the main growth panels. A single
invalid source as $K$ increases is a separate dilution diagnostic. Use a
declared coupling of site populations and target samples across $K$ so that
changes in case mix do not masquerade as source-count effects.
\item \textbf{Compare under matched information.} Include calibrated target-only,
current RoCE, equally informed known-valid references, oracle-valid-set
diagnostics, and appropriate Guo/RIFL adaptations. Report coverage, expected
length, point RMSE, fallback frequency, budget size and runtime. Deliberately
wrong valid-count assumptions and active clipping are boundary cases.
Distinguish oracles for valid-source identities, nuisance functions and
covariance; their information is not interchangeable.
\end{enumerate}

The earlier 124-setting Gaussian study and 90 fitted-score cases inform these
design choices; they do not validate the procedure defined here. No new
large-scale fitting campaign is part of this design document. The principal
methodological question is whether the explicit error protection can be made
small enough to retain useful borrowing at realistic $n,p,K$.

\paragraph{Checks accompanying this design version.}
The matrix identities, variance domination and bias geometry were checked in
60 deterministic-seed matrix cases. An independent enumeration of all
262,144 configurations in a small randomized binary-outcome model checked
the patient-pair reconstruction, exact expectation and variance, and the
one-sided inversion. The example has a nonconstant shared target prediction,
within-person arm dependence and different valid source arms. There are
696 passing assertions, including inversion boundaries, the zero-dispersion
case, the Bernstein radius and covariance-bound matrix checks. These are
mathematical implementation checks, not a
fitted-method coverage or efficiency experiment.

\paragraph{Implementation contract.}
The research interface should expose arguments for the minimum valid-source
counts (\texttt{valid\_source\_min}), valid-score bias budgets
(\texttt{valid\_score\_bias}), covariance certificate
(\texttt{covariance\_certificate}), interval calibration
(\texttt{interval\_calibration}), split seed (\texttt{split\_seed})
and nuisance solver (\texttt{nuisance\_solver}).
The first design uses \texttt{common\_outcome} scores,
\texttt{honest\_pilot} sample roles and \texttt{second\_moment} calibration.
Missing bias budgets or covariance certificates must not silently become
zero-bias or oracle-covariance assumptions. These are proposed interfaces,
not claims that a production implementation already exists.

\appendix
\section{An explicit covariance upper bound}\label{app:covariance}

Conditional on training, recenter each patient-score coordinate by a known
constant so that $|G_{b,i,l}|\le L_{b,l}$ on its full support. Centering does
not change covariance. Bounds may come from bounded outcomes and the actual
weight caps; a sample maximum is not a population bound. Known constant
coordinates can have $L_{b,l}=0$.

For pilot size $m_b$, write $\widehat\mu_l=\mathbb P_{b,\calP}G_l$ and
$\widehat M_{lk}=\mathbb P_{b,\calP}G_lG_k$. Allocate $\delta_b$ across the
$d_b+d_b(d_b+1)/2=N_b$ mean and second-moment entries. Hoeffding gives
simultaneously, with probability at least $1-\delta_b$,
\[
 |\widehat\mu_l-\mu_l|\le r_l=L_l\sqrt{2\log(2N_b/\delta_b)/m_b},\qquad
 |\widehat M_{lk}-M_{lk}|\le r_{lk}=L_lL_k\sqrt{2\log(2N_b/\delta_b)/m_b}.
\]
Let $\widehat\Omega_b^P=\widehat M-\widehat\mu\widehat\mu^T$ and
\[
 E_{lk}=r_{lk}+|\widehat\mu_l|r_k+|\widehat\mu_k|r_l+r_lr_k,
 \quad e_{b,l}=\sum_k E_{lk},\quad
 U_b^{\rm mom}=\widehat\Omega_b^P+\diag(e_{b,1},\ldots,e_{b,d_b}).
\]
Symmetric diagonal dominance gives $\Omega_b\preceq U_b^{\rm mom}$ on
that event. A second, deterministic upper bound is
$U_b^{\rm range}=d_b\diag(L_{b,1}^2,\ldots,L_{b,d_b}^2)$, since
$\Var(v^TG_b)\le(\sum_l|v_l|L_{b,l})^2\le v^TU_b^{\rm range}v$.
Choose $\overline\Omega_b$ as the one of these two matrices with smaller
trace, using a fixed tie rule. Both are upper bounds on the confidence event,
so this pilot-only selection preserves coverage; it is not an entrywise
minimum. Row-specific inflation preserves known constant coordinates.
Set
$W=L\diag\{\overline\Omega_b/n_b^E\}_bL^T$ and
$\sum_b\delta_b\le\delta_\Sigma$. Check positive definiteness before solving.
This is fully computable given valid range bounds, but can be extremely loose
when weight caps are large. Boundedness is needed here even though
Proposition~\ref{prop:ustat} itself only needs second moments. Replacing this
certificate by an unproved empirical covariance would change the guarantee.

\section{Proof of the quadratic reduction}\label{app:quadratic}

Put $A=L^TD_WL$. Within a site,
\[
 \overline G_b^TA_{bb}\overline G_b-
 \tr\{A_{bb}S_b^E/n_b^E\}
 =\frac{\sum_{i\ne k}G_{bi}^TA_{bb}G_{bk}}{n_b^E(n_b^E-1)}.
\]
Cross-site terms are the ordinary products of independent means. This gives
a quadratic form whose entire within-patient diagonal blocks are zero, and
establishes unbiasedness. After centering each patient block, its linear
part has variance $4r^TD_W\Sigma_ED_Wr$. Its quadratic part has variance
\begin{align*}
 &2\sum_{b\ne c}\frac{\tr(A_{bc}\Omega_cA_{cb}\Omega_b)}{n_b^En_c^E}
 +2\sum_b\frac{\tr\{(A_{bb}\Omega_b)^2\}}{n_b^E(n_b^E-1)}\\
 &\hspace{2em}\le2\kappa_E\tr\{(D_W\Sigma_E)^2\}.
\end{align*}
The linear and quadratic parts are uncorrelated: each putative third-moment
term contains an independent centered patient block occurring only once.
No within-patient independence or symmetry is needed. Since
$D_WWD_W=D_W$ and $\Sigma_E\preceq W$,
\[
 r^TD_W\Sigma_ED_Wr\le\calB,
 \qquad \tr\{(D_W\Sigma_E)^2\}\le\nu.
\]
This proves~\eqref{eq:energy-var}. Cantelli bounds
$\calB-\widehat{\calB}$ by
$\sqrt{c_D(4\calB+2\kappa_E\nu)}$ outside an event of probability
$\alpha_D$. Solving the resulting quadratic inequality gives
\eqref{eq:energy-upper}. A negative discriminant has no feasible positive
energy on the concentration event; the stated nonnegative convention is safe.

\section{Bias geometry and a count-based upper bound}\label{app:count-bound}

The identities $P_WWP_W=P_W$, $R_{\calA}WR_{\calA}=R_{\calA}$ and
$P_WWR_{\calA}=R_{\calA}$ follow from $H^TR_{\calA}=0$ and the GLS definitions.
After conjugation by $W^{1/2}$ these are nested orthogonal projectors.
Their difference is therefore positive semidefinite, has rank
$2(K+1)-|\calA|$, and satisfies $D_WWD_W=D_W$.

For $H^Td_V=\ell$, direct multiplication gives
\[
 \widetilde d_V-a_W=D_WWd_V,\qquad
 C_V=d_V^TWD_WWd_V\ge0.
\]
Consequently $|(\widetilde d_V-a_W)^Tr|\le\sqrt{C_V\calB}$.
The reference projection adds coefficients only on $\calA$, so
$\widetilde d_V$ remains supported on the true valid set; this proves
\eqref{eq:bias-geometry}. In particular, residual disagreement alone cannot
control a common drift. The nonzero budgets in $B_V$ must remain.

For the shared-target block structure, replace each private source block
\[
 \begin{pmatrix}d_{1j}&c_j\\c_j&d_{0j}\end{pmatrix}
 \quad\text{by}\quad
 \begin{pmatrix}d_{1j}+|c_j|&0\\0&d_{0j}+|c_j|\end{pmatrix}.
\]
The difference is positive semidefinite; call the resulting matrix
$\overline W\succeq W$. This is an explicit conservative bound, not an
assertion that private arm covariance is zero. For each admissible $V$, take
$d_V$ as GLS under $\overline W_V$. Then
$C_V\le v_V(\overline W)-v_W$.

Writing the target anchor/common covariance blocks as $A_0,B_0,C_0$, the
selected private pool variances $u_a$ give source-pool contrast coefficients
\[
 \beta(u_1,u_0)=\{A_0+C_0-B_0-B_0^T+\diag(u_1,u_0)\}^{-1}
                    (A_0-B_0^T)\ell.
\]
Anchor coefficients are $\ell-\beta$ and within-arm source coefficients use
positive inverse-private-variance shares. Selecting the largest private
variances in each arm maximizes $v_V(\overline W)$ for the prescribed counts,
including any mandatory known-valid coordinates. The smallest and largest
allowed pools bound $(u_1,u_0)$ in a rectangle.

Set $\widetilde b=|I-R_{\calA}W|^Tb$. A conservative common-family bias envelope
is the maximum of
\[
 \sum_a\left\{|\ell_a-\beta_a|\widetilde b_{t,a}
       +|\beta_a|\max_j\widetilde b_{j,a}\right\}
\]
over that rectangle. With positive private variances and a positive definite
displayed matrix throughout it, each coefficient follows a common
linear-fractional parameter along one coordinate. Convexity of the weighted
absolute-value sum makes the four corners sufficient. An arm with no
guaranteed sources fixes its pool coefficient to zero and uses the remaining
one-dimensional calculation. The resulting pair
$(\overline B_{\rm count},\overline C_{\rm count})$ supplies
\eqref{eq:bias-radius}. Exact subset calculations must verify this shortcut on
small instances before it is used for large $K$.

\bibliographystyle{plainnat}
\bibliography{references}
\end{document}
